% ==============================================================================
% SEÇÃO 1: DESCRITIVO DETALHADO DO SISTEMA VITALIS v2.0
% ==============================================================================

\section{Descritivo Detalhado do Sistema Vitalis v2.0}
\label{sec:descritivo}

O Sistema de Gestão Hospitalar Inteligente Vitalis (Versão 2.0) é uma plataforma integrada de prontuário eletrônico do paciente (PEP), regulação de leitos e suporte à decisão clínica, projetada para atender às demandas de hospitais gerais e de alta complexidade. A solução combina a gestão dos processos administrativos e assistenciais com modelos de inteligência artificial preditiva e generativa, atuando como camada consultiva para médicos, enfermeiros, farmacêuticos e gestores. Esta seção apresenta o contexto operacional da instituição hospitalar, a arquitetura de microsserviços, o ciclo assistencial completo do paciente, os módulos especializados de inteligência artificial e as salvaguardas éticas e legais do projeto.

\subsection{Contexto Operacional e Visão de Negócio}
\label{subsec:contexto_negocio}

Ambientes hospitalares terciários operam sob contínua sobrecarga cognitiva, com fluxos assistenciais simultâneos em unidades de pronto atendimento, enfermarias de internação, centros cirúrgicos e unidades de terapia intensiva (UTI). Nesses cenários, a fragmentação dos dados entre sistemas legados desconectados gera gargalos crônicos:

\begin{itemize}
    \item Retardo na identificação de pacientes em risco de deterioração clínica rápida ou sepse nas primeiras horas de triagem;
    \item Erros decorrentes da transcrição manual repetitiva de sinais vitais, queixas clínicas e planos terapêuticos;
    \item Ociosidade invisível de leitos em decorrência da falta de sincronismo entre a alta médica, o serviço de higienização predial e a regulação central;
    \item Glosas financeiras frequentes provocadas por inconsistências documentais entre as prescrições médicas executadas e os itens faturados nas guias da saúde suplementar.
\end{itemize}

O Vitalis v2.0 unifica esses domínios em uma plataforma digital integrada. O propósito operacional é reduzir o tempo de espera no acolhimento, maximizar a rotatividade segura dos leitos por meio da antecipação da demanda assistencial e zerar glosas administrativas decorrentes de falhas de preenchimento. A plataforma prioriza a segurança do paciente e o rigor assistencial, de modo que todo modelo preditivo ou generativo atua estritamente como suporte técnico subordinado à homologação do profissional responsável \cite{topol:2019}.

\subsection{Arquitetura em Microsserviços e Engenharia de Software}
\label{subsec:arquitetura_software}

Para garantir disponibilidade contínua, tolerância a falhas e isolamento de domínios funcionais, o Vitalis v2.0 adota uma arquitetura orientada a microsserviços fundamentada nas seguintes escolhas de engenharia:

\begin{itemize}
    \item \textbf{Camada de Serviços e APIs}: Cada microsserviço é construído sobre o framework NestJS utilizando o motor de rede Fastify, o que assegura baixa latência sob alta concorrência de conexões. O código adota TypeScript com separação em camadas inspirada na Arquitetura Limpa, isolando entidades de domínio, casos de uso e adaptadores de infraestrutura.
    \item \textbf{Orquestração de Filas e Tarefas Assíncronas}: O sistema emprega o BullMQ suportado por instâncias em memória do Redis. Essa estrutura gerencia o processamento assíncrono de tarefas com alto custo computacional, tais como chamadas a modelos de linguagem, indexação vetorial de literatura médica, geração de laudos em PDF e compilação de lotes de guias TISS. Filas prioritárias garantem que alertas de deterioração clínica sejam despachados em milissegundos, enquanto rotinas financeiras operam em baixa prioridade.
    \item \textbf{Persistência Relacional em Terceira Forma Normal (3NF)}: A persistência operacional utiliza o PostgreSQL. O esquema de dados é normalizado rigorosamente em 3NF \cite{pressman:2016}, eliminando redundâncias indesejadas e inconsistências funcionais. O modelo estrutura tabelas nucleares como \texttt{PACIENTE}, \texttt{ATENDIMENTO}, \texttt{PRONTUARIO\_LOGICO}, \texttt{EVOLUCAO}, \texttt{PRESCRICAO}, \texttt{LEITO} e \texttt{GUIA\_TISS}, utilizando índices compostos otimizados para consultas de séries temporais de sinais vitais e histórico do prontuário.
    \item \textbf{Controle de Acesso Granular (RBAC e ABAC)}: A autorização combina papéis hierárquicos com atributos de contexto, tais como registro ativo em conselhos profissionais (CRM, COREN, CRF), especialidade clínica cadastrada e vinculação formal ao episódio de cuidado ativo do paciente. Usuários administrativos não possuem permissão de acesso a notas confidenciais de saúde mental.
    \item \textbf{Comunicação em Tempo Real}: A sincronização do mapa de leitos e o disparo de alertas emergenciais utilizam WebSockets bidirecionais, atualizando imediatamente os terminais dos postos de enfermagem sem exigir recarregamento manual de página.
    \item \textbf{Gateway de IA Seguro e Anonimizador}: Todas as solicitações direcionadas a serviços de IA transitam obrigatoriamente por um gateway intermediário. Esse serviço executa inspeção de payload, remove identificadores pessoais e formata os dados clínicos antes da inferência, protegendo a privacidade dos dados.
    \item \textbf{Resiliência via Circuit Breaker}: Se um serviço externo de inferência atingir latência superior a 2,0 segundos ou apresentar falha de resposta, o sistema chaveia automaticamente para mecanismos determinísticos locais baseados em árvores de decisão e tabelas clínicas tradicionais, preservando a continuidade operacional sem interrupções.
\end{itemize}

\subsection{Ciclo Assistencial Completo do Paciente}
\label{subsec:ciclo_assistencial}

O Vitalis v2.0 gerencia todo o percurso assistencial e administrativo do paciente no hospital e assegura a rastreabilidade integral em cada etapa do fluxo.

\subsubsection{Recepção e Admissão Cadastral}
O ciclo se inicia com a identificação do paciente na recepção hospitalar. O sistema valida os dígitos verificadores do CPF junto às regras cadastrais da Receita Federal e consulta o número do Cartão Nacional de Saúde (CNS). Em casos de emergência com pacientes inconscientes ou não identificados, o sistema permite a abertura imediata de prontuário sob o protocolo provisório de "Paciente Desconhecido", exigindo regularização cadastral posterior. Para pacientes de saúde suplementar, o módulo executa a verificação eletrônica de elegibilidade da operadora de saúde e gera o token de atendimento, posicionando o paciente na fila de acolhimento.

\subsubsection{Acolhimento e Triagem Aumentada por IA}
Na sala de classificação de risco, o enfermeiro triador realiza a avaliação do paciente fundamentada no Protocolo de Manchester \cite{manchester:2014}. O sistema disponibiliza entrada via formulário estruturado e módulo de reconhecimento de voz (Speech-to-Text via Whisper), transcrevendo as queixas do paciente e convertendo a linguagem natural em descritores clínicos objetivos.

Enquanto o profissional afere os parâmetros vitais (frequência cardíaca, pressão arterial sistólica e diastólica, saturação periférica de oxigênio, frequência respiratória e escala de coma de Glasgow), o Vitalis v2.0 executa em paralelo o cálculo automatizado de dois escores clínicos de gravidade:
\begin{itemize}
    \item \textbf{qSOFA} (\textit{Quick Sequential Organ Failure Assessment}): Pontua alterações em frequência respiratória ($\ge 22\,\text{rpm}$), pressão arterial sistólica ($\le 100\,\text{mmHg}$) e estado neurológico, alertando para risco elevado de sepse em pacientes fora da UTI \cite{singer_sepsis:2016};
    \item \textbf{NEWS2} (\textit{National Early Warning Score 2}): Agrega os parâmetros fisiológicos para mensurar o grau de estresse hemodinâmico e prever risco iminente de colapso respiratório ou necessidade de leito intensivo.
\end{itemize}
Se a classificação resultar na prioridade Vermelha (Emergência), o sistema ativa avisos visuais e sonoros nos monitores da sala de emergência e transfere o paciente para o topo absoluto da fila médica.

\subsubsection{Regulação e Alocação Preditiva de Leitos}
Após a avaliação médica indicar a necessidade de internação, entra em ação o módulo de gestão de leitos. A plataforma apresenta um mapa visual com o estado de cada vaga: Livre, Reservado, Ocupado, Em Higienização, Manutenção ou Isolamento.

O motor analítico de previsão de tempo de permanência (\textit{Length of Stay}, LOS) utiliza algoritmos treinados sobre variáveis demográficas, diagnóstico principal CID-10, histórico de comorbidades e complexidade de procedimentos solicitados. O sistema estima o período médio de internação e indica a melhor acomodação disponível por ala e especialidade. O fluxo proíbe a transferência para leitos marcados como "Em Higienização" antes da validação digital pela equipe de hotelaria hospitalar.

\subsubsection{Prontuário Eletrônico (PEP) e Copiloto Farmacoterapêutico}
Durante a estada hospitalar, todas as anotações clínicas, ordens de serviço, solicitações de exames diagnósticos e prescrições médicas são concentradas no PEP. O módulo de prescrição eletrônica conta com um copiloto farmacoterapêutico baseado em arquitetura RAG (\textit{Retrieval-Augmented Generation}) \cite{lewis_rag:2020}.

Esse copiloto recupera informações em uma base vetorial que reúne bulários, guias de diluição, monografias de interações e diretrizes clínicas institucionais. Ao selecionar um medicamento, o motor faz a checagem cruzada das seguintes variáveis:
\begin{itemize}
    \item Alergias ativas e reações adversas prévias registradas no perfil do paciente;
    \item Interações farmacológicas entre os itens da prescrição atual e prescrições vigentes;
    \item Ajuste de dosagem por faixa de peso corporal e função renal calculada com base na creatinina sérica recente;
    \item Duplicidade de classes terapêuticas e incompatibilidades físico-químicas de administração pela mesma via venosa.
\end{itemize}
Caso seja detectada uma incompatibilidade moderada ou grave, o sistema exibe alertas contextualizados e exige a concordância justificada do médico para autorizar o despacho à farmácia satélite.

\subsubsection{Evoluções Clínicas e Sumário de Alta Automatizado}
Os médicos e a equipe multiprofissional registram evoluções diárias estruturadas por sistemas orgânicos (Cardiovascular, Respiratório, Neurológico, Gastrointestinal e Renal). Quando o paciente atinge estabilidade clínica e preenche os critérios para desospitalização, o médico aciona a geração automatizada do sumário de alta.

O modelo de linguagem agrega os marcos temporais da internação: motivo da admissão, diagnósticos confirmados, procedimentos cirúrgicos executados, exames laboratoriais e radiológicos com resultados anormais, fármacos administrados e resposta terapêutica. O sistema compõe uma minuta formal contendo o resumo do internamento, a receita médica de continuidade domiciliar e o plano de retorno ambulatorial. O médico analisa o texto, realiza as modificações pertinentes e assina o documento.

\subsubsection{Faturamento Hospitalar e Auditoria TISS/TUSS}
Após a assinatura da alta, o atendimento é disponibilizado ao setor de faturamento. O sistema compila automaticamente os itens consumidos (diárias de quarto, taxas de uso de sala cirúrgica, gases medicinais, exames realizados, medicamentos dispensados e honorários profissionais).

O motor de auditoria preventiva confronta esses lançamentos com a tabela de cobertura do convênio e as regras de codificação TUSS estabelecidas pela Agência Nacional de Saúde Suplementar (ANS) \cite{ans_tiss:2023}. O sistema detecta preventivamente discrepâncias como cobrança de antibióticos de uso restrito sem laudo de cultura comprobatório anexado, cobranças de taxas cirúrgicas incompatíveis com o tempo de anestesia ou ausência de autorização prévia da operadora. Isso permite a correção documental pelo faturista antes da transmissão do lote XML, reduzindo a incidência de glosas.

\subsection{Módulos Especializados de Inteligência Artificial e Suporte à Decisão}
\label{subsec:modulos_ia}

A camada de inteligência do Vitalis v2.0 é organizada em quatro pipelines de processamento focados em pontos críticos do atendimento:

\begin{itemize}
    \item \textbf{Processamento de Linguagem Natural e Reconhecimento de Voz}: Responsável por converter o áudio da anamnese em texto estruturado durante a triagem. O módulo utiliza reconhecimento fonético calibrado para terminologia médica em língua portuguesa, identificando entidades clínicas como sintomas, dosagens e tempo de evolução da queixa.
    \item \textbf{Motor Preditivo de Risco Hemodinâmico}: Calcula continuamente os escores qSOFA e NEWS2 a partir de novas entradas de sinais vitais e exames laboratoriais. O modelo opera com pesos calibrados para evitar alarmes espúrios, mantendo alta sensibilidade para detecção precoce de sepse e insuficiência respiratória aguda.
    \item \textbf{Recuperação Vetorial Farmacoterapêutica (RAG)}: Utiliza embeddings densos armazenados em banco vetorial para correlacionar o histórico clínico do paciente com a literatura de segurança farmacêutica. O sistema recupera trechos específicos de bulários e diretrizes da Anvisa para embasar os alertas exibidos ao prescritor.
    \item \textbf{Síntese Clínica e Mitigação de Alucinações}: O motor gerador de sumários de alta emprega técnicas de engenharia de prompt com restrição estrita de contexto (\textit{grounding}). O modelo processa exclusivamente os dados presentes no prontuário eletrônico do paciente, sem inferir informações ausentes e sinalizando lacunas documentais ao médico revisor.
\end{itemize}

\subsection{Governança Clínica, Segurança da Informação e Conformidade Regulatória}
\label{subsec:governanca_conformidade}

A aplicação de tecnologias de inteligência artificial em ambientes hospitalares impõe rígidos limites operacionais e conformidade com os marcos legais vigentes no Brasil:

\begin{itemize}
    \item \textbf{Princípio do Human-in-the-Loop}: O sistema estabelece que nenhuma inferência de IA possui caráter deliberativo autônomo. Diagnósticos, cálculos de risco de sepse, recomendações posológicas e minutas de documentos clínicos funcionam estritamente como ferramentas de auxílio cognitivo. Toda conduta exige a avaliação, confirmação e assinatura expressa do profissional assistente legalmente habilitado.
    \item \textbf{Conformidade com o Conselho Federal de Medicina (CFM)}: O Vitalis v2.0 observa integralmente as diretrizes da Resolução CFM nº 1.821/2007 \cite{cfm_prontuario:2007}, que rege o uso de sistemas informatizados para a guarda e manuseio de prontuários, e da Resolução CFM nº 2.299/2021 \cite{cfm_digital:2021}, que disciplina a emissão de documentos médicos eletrônicos. Toda evolução clínica, prescrição e sumário de alta é assinado digitalmente com certificados no padrão ICP-Brasil (PAdES com carimbo de tempo). Prontuários encerrados tornam-se imutáveis contra alterações diretas, admitindo apenas a inclusão de notas aditivas de retificação auditadas.
    \item \textbf{Conformidade com a Lei Geral de Proteção de Dados (LGPD)}: O tratamento de dados pessoais sensíveis de saúde é conduzido sob as bases legais da tutela da saúde e da proteção à vida, preconizadas no Artigo 7º, inciso VIII, e no Artigo 11, inciso II, alínea "f", da Lei nº 13.709/2018 \cite{brasil_lgpd:2018}. Antes de qualquer transmissão para APIs de IA, o pipeline de sanitização remove os 18 identificadores diretos de dados pessoais (PII), incluindo nomes, números de documento, endereços e identificadores hospitalares, operando sobre dados pseudonimizados.
    \item \textbf{Interoperabilidade com HL7 FHIR}: O barramento de integração implementa o padrão HL7 FHIR Release 4 \cite{hl7_fhir:2019}, permitindo a comunicação com sistemas de laudos externos (PACS/LIS), operadoras de planos de saúde e redes do Sistema Único de Saúde (RNDS).
    \item \textbf{Trilha de Auditoria e Explicabilidade}: A tabela \texttt{LOG\_INFERENCIA\_IA} armazena o histórico auditável de cada chamada aos motores de inteligência artificial. O registro inclui o identificador do atendimento, modelo e versão utilizados, hash criptográfico do prompt enviado, tempo de resposta em milissegundos, nível de confiança estatística, sugestão emitida e a decisão do profissional (aprovada, rejeitada ou alterada, acompanhada da devida justificativa técnica).
    \item \textbf{Disponibilidade e Recuperação de Desastres}: Para garantir a continuidade operacional em um serviço essencial de saúde, a infraestrutura define objetivos estritos de recuperação: RTO (\textit{Recovery Time Objective}) de até 4 horas e RPO (\textit{Recovery Point Objective}) de até 1 hora, sustentados por réplicas síncronas de banco de dados e rotinas diárias de backup distribuídas geograficamente.
\end{itemize}
